\scanpage{58}
% E: text|upper left date|3.1.84
3.1.84
\deleted{
% E: text|upper left title three lines|Transport parallèle le long
Transport parallèle le long d'un demi-tour, \deleted{\unsure{au le|au}} le long d'une portion ultra \deleted{\unsure{singulière de polyg.}} \inserted{supersingulière}
% E: display|upper left boxed formula|lambda(t)
\[\boxed{\lambda(t)=t-(n-1)}\]
% E: text|left upper middle|NB Le pgcd
\uline{NB} Le pgcd de $n-1$ et de $2n$ est $1$ si $n$ pair, $2$ si $n$ impair.
% E: text|left middle introductory sentence|Donc on voit que l'image
Donc on voit que l'image de
% E: diagram|left middle|pi1(X) to D2n
\diagram{\begin{tikzcd}
\pi_1(X)\arrow[r] & \deleted{\unsure{\mathbb Z}}\mathbb D_{2n}
\end{tikzcd}}
% E: text|left middle image and index|a une image qui contient
 a une image qui contient $\mathbb D_{2n}^{+}\inserted{\unsure{{}\simeq\mathbb Z/2n\mathbb Z}}$ si $n$ pair (donc cette image est d'indice $1$ ou $2$) \deleted{ou bien elle} \inserted{et elle contient $\mathbb Z/n\mathbb Z$ si $n$ impair} (donc elle est tjrs d'indice $1,2$ ou $4$).
% E: text|left lower middle cutting and translation|Si on prend le découpage
\deleted{\unsure{Cela ne serait vrai}} Si on prend le \unsure{découpage} $\widetilde X$ de $X$ suivant les composantes supersingulières de $K$, on trouve dans l'image de
% E: diagram|left lower|pi1(X tilde) to D2n
\diagram{\begin{tikzcd}
\pi_1(\widetilde X)\arrow[r] & \mathbb D_{2n}
\end{tikzcd}}
% E: text|left lower|la translation
la translation
% E: display|left bottom|t to lambda prime(t)
\[t\longrightarrow\lambda'(t)=t-2(n-1).\]
% E: text|left bottom final two lines|Or le pgcd
Or le pgcd de $2(n-1)$ et \unsure{$n$|$2n$} est \unsure{un div. tjrs $2$|ici tjrs $2$}.
Donc l'image contient tjrs $\mathbb Z/n\mathbb Z$, d'indice $4$
% E: text|right top continuation|le quotient étant D2
(le quotient étant $\mathbb D_2\simeq\mathbb F_2\times\mathbb F_2$)
% E: text|right upper paragraph|Ceci ne nous dit rien
Ceci ne nous dit \inserted{rien} sur l'image de $\pi_1(X)$, $\pi_1(\widetilde X)$ dans $\mathbb D_2$, et en particulier sur son image dans $\pm1$, \inserted{\unsure{par}} composée avec
% E: diagram|right upper middle|D2n to plus or minus 1
\diagram{\begin{tikzcd}
\mathbb D_{2n}\arrow[r,"\chi"] & \pm1
\end{tikzcd}}
% E: text|right middle and lower|On a l'impression
On a l'impression que pour avoir des informations, il faut faire pivoter une \unsure{p.r.|por.} autour d'un sommet de $\uline{\Sigma}$ et faire ainsi un cycle et voir l'effet produit. On va faire un demi-tour, à partir d'une position supersingulière disons, par transport parallèle l'orientation de la position supersingulière aura changé, et \deleted{$\unsure{\chi}$} $\chi(\lambda)$ sera égal : $+1$ si le nb total des ``pas'' pour revenir au pt de départ est \uline{impair}; $-1$ dans le cas contraire. Donc on devrait pouvoir le faire en un nb \uline{pair} de pas !
}
% E: margin note|diagonal upper middle|Toute la partie
\margin{\unsure{Toute la partie combinée !|Toute la partie combinatoire !}}
% E: margin note|diagonal centre lower dated 4.1|Attention
\margin{4.1. Attention ! Ces \unsure{positions|portions} sont exclues dans le jeu des $\operatorname{Pol}(\unsure{D|CD})$}
% E: margin note|diagonal bottom centre|cf notes du 4.1
\margin{\unsure{Génial|Général} cf notes du 4.1.}
\mspage{57}
